\documentclass[12pt,a4paper]{article}

% --- Paquetes básicos ---
\usepackage[spanish]{babel}     % Español
\usepackage[utf8]{inputenc}     % Acentos
\usepackage[T1]{fontenc}
\usepackage[spanish]{babel}     % Idioma
\usepackage{geometry}           % Márgenes
\usepackage{graphicx}           % Imágenes
\usepackage{float}              % Posición exacta de figuras
\usepackage{hyperref}           % Hipervínculos
\usepackage{xcolor}             % Colores
\usepackage{fancyhdr}           % Encabezados y pies de página
\usepackage{titlesec}           % Formato de títulos
\usepackage{listingsutf8}       % Soporte UTF-8 en listing
\usepackage{listings}           % Código
\usepackage{caption}            % Personalizar captions
\usepackage{amsmath, amssymb}   % Símbolos matemáticos
\usepackage{booktabs}           % Tablas elegantes
\usepackage{listings}           % Código
\usepackage{longtable}
\usepackage[table]{xcolor}
\usepackage{listings}
\usepackage{xcolor}
\usepackage{seqsplit}


% --- Configuración de códigos ---
\lstset{
  basicstyle=\ttfamily\small,
  backgroundcolor=\color{gray!10},
  frame=single,
  breaklines=true,
  columns=fullflexible,
  inputencoding=utf8,
  extendedchars=true,
  literate=
    {á}{{\'a}}1 {é}{{\'e}}1 {í}{{\'i}}1 {ó}{{\'o}}1 {ú}{{\'u}}1
    {Á}{{\'A}}1 {É}{{\'E}}1 {Í}{{\'I}}1 {Ó}{{\'O}}1 {Ú}{{\'U}}1
    {ñ}{{\~n}}1 {Ñ}{{\~N}}1
}

% Estilo para Bash
\lstdefinestyle{miEstiloBash}{
    language=bash,
    basicstyle=\small\ttfamily,
    backgroundcolor=\color{gray!10},
    keywordstyle=\color{blue}\bfseries,
    commentstyle=\color{green!60!black},
    stringstyle=\color{red},
    showstringspaces=false,
    frame=single,
    breaklines=true,
    captionpos=b
}

% Estilo para Scala
\lstdefinestyle{miEstiloScala}{
    language=Scala,
    basicstyle=\small\ttfamily,
    backgroundcolor=\color{gray!5},
    keywordstyle=\color{blue}\bfseries,
    commentstyle=\color{green!60!black}\itshape,
    stringstyle=\color{red},
    showstringspaces=false,
    frame=single,
    breaklines=true,
    captionpos=b,
    morekeywords={val, var, def, class, object, trait, extends, with, match, case, new, if, else, for, while, return, type, import, package}
}

% Estilo para Java
\lstdefinestyle{miEstiloJava}{
    language=Java,
    basicstyle=\small\ttfamily,
    backgroundcolor=\color{gray!5},
    keywordstyle=\color{blue}\bfseries,
    commentstyle=\color{green!60!black}\itshape,
    stringstyle=\color{red},
    showstringspaces=false,
    frame=single,
    breaklines=true,
    captionpos=b
}

% --- Configuración de márgenes ---
\geometry{top=2.5cm, bottom=2.5cm, left=3cm, right=2.5cm}

% --- Configuración de hipervínculos ---
\hypersetup{
    colorlinks=true,
    linkcolor=blue,
    urlcolor=cyan
}

% --- Encabezados y pies de página ---
\pagestyle{fancy}
\fancyhf{}
\fancyhead[L]{Universidad de Valladolid}
\fancyhead[R]{Máster en Ingeniería Informática}
\fancyfoot[C]{\thepage}

% --- Formato de títulos ---
\titleformat{\chapter}[hang]{\bfseries\huge}{\thechapter.}{1em}{}
\titleformat{\section}[hang]{\bfseries\Large}{\thesection}{1em}{}

% -- Información para la portada
\title{\LARGE \textbf{Proyecto de Clasificación: Primera Etapa}}
\author{{\large \textbf{Autores:}}\\[0.2cm]
    Gonzalo Agúndez Sáez \\
    Paula Negro Asegurado \\
    Diego Sanz Santos }
\date{\normalsize \today}

\makeatletter
\let\thetitle\@title
\let\theauthor\@author
\let\thedate\@date
\makeatother

\begin{document}

\begin{titlepage}
    \centering
    \vspace*{1.5cm}
    \includegraphics[width=0.7\textwidth]{logoInfo.png}\\[1cm] % Descomentar si se tiene el logo

    {\Huge \bfseries \thetitle}\\[1.5cm]

    {\Large \bfseries Asignatura: Técnicas escalables de análisis de datos (Clasificadores)}\\[0.3cm]
    {\large Escuela de Ingeniería Informática}\\[1.5cm]

    \theauthor \\[0.3cm]
    \textbf{Curso:} 2026 -- 2027 \\[0.3cm]
    \textbf{Fecha:} \thedate

    \vfill
    \rule{0.5\linewidth}{0.5pt}\\
    \textit{Universidad de Valladolid}
\end{titlepage}

\tableofcontents
\newpage
\section{Introducción}
% Breve introducción al documento

\section{Resumen Ejecutivo del Conjunto de Datos}
\begin{itemize}
    \item \textbf{Origen de los datos:} competición de Kaggle \textit{New York City Taxi Trip Duration}, basada en los registros de taxis amarillos de Nueva York de 2016 publicados por la NYC Taxi and Limousine Commission (TLC). Se ha utilizado el fichero \texttt{train.csv} disponible en \url{https://www.kaggle.com/c/nyc-taxi-trip-duration/data}.
    \item \textbf{Propósito:} El proposito es predecir la duración de un viaje en taxi a partir de sus atributos (tiempo, coordenadas, pasajeros, etc.). En este trabajo se emplean para un problema de \textbf{clasificación multiclase}: predecir si un viaje es corto, medio o largo.
    \item \textbf{Tamaño:} 1\,458\,644 instancias y 11 atributos originales.
    \item \textbf{Atributos (11 en total):}
    \begin{itemize}
        \item \textbf{Nominales (4):} \texttt{id} (identificador único del viaje), \texttt{vendor\_id} (proveedor del registro) y las marcas temporales \texttt{pickup\_datetime} y \texttt{dropoff\_datetime}.
        \item \textbf{Numéricos continuos (4):} \texttt{pickup\_longitude}, \texttt{pickup\_latitude}, \seqsplit{\texttt{dropoff\_longitude}} y \texttt{dropoff\_latitude}.
        \item \textbf{Numéricos discretos (2):} \texttt{passenger\_count} y \texttt{trip\_duration} (en segundos).
        \item \textbf{Booleanos (1):} \texttt{store\_and\_fwd\_flag} (Y/N).
    \end{itemize}
    \item \textbf{Descripción de la clase:} la clase se construye a partir de \texttt{trip\_duration}, que se descarta como atributo predictor. Se definen tres etiquetas: \textit{short} (duración $\leq$ 10 min), \textit{medium} (entre 10 y 25 min) y \textit{long} ($>$ 25 min). La distribución es: \textit{short} 44,76\,\%, \textit{medium} 42,87\,\%, \textit{long} 12,37\,\%.
\end{itemize}

%Limpieza inicial ("tonta"): Eliminar los registros que estén vacíos o sean no válidos en el DataFrame original.   División: Crear las particiones para los conjuntos de entrenamiento y prueba.   Limpieza profunda, selección y transformación:Tratar los valores ausentes, inconsistentes, el ruido y los outliers.   Seleccionar atributos y construir nuevos a partir de los originales.   Realizar las transformaciones necesarias para obtener las columnas features y labels (aquí es donde transformarás la duración a "short", "medium" y "long").   La regla de oro en el tercer paso es que cualquier ajuste de un modelo de transformación (estimator) se entrena o ajusta (hacer fit) únicamente sobre el conjunto de entrenamiento, y luego se aplica exactamente igual al conjunto de prueba.%

\section{Exploración de Datos}

\subsection{Descripción del problema}

El problema consiste en predecir cuánto va a durar un viaje en taxi amarillo en
Nueva York a partir de datos como la hora y el lugar de recogida, el destino, el
número de pasajeros y el proveedor. El planteamiento original es de regresión,
pero nosotros lo tratamos como un problema de \textbf{clasificación multiclase}:
a partir de \texttt{trip\_duration} (en segundos) definimos tres clases,
\textit{short} (hasta 10 minutos), \textit{medium} (entre 10 y 25) y
\textit{long} (más de 25). Un modelo así podría servir, por ejemplo, para avisar
al cliente de cuánto tardará su trayecto o para organizar mejor la flota.

Los datos son viajes de taxis amarillos de 2016 publicados por la TLC de Nueva
York. Kaggle los muestreó y depuró para la competición \textit{New York City Taxi
Trip Duration}, un reto abierto pensado para compartir código y datos y aprender
entre todos. Los viajes del conjunto de entrenamiento empezaron entre el 1 de
enero y el 30 de junio de 2016, mientras que la competición se celebró entre
julio y septiembre de 2017.

En la competición el objetivo era predecir la duración exacta de cada viaje
(regresión, evaluada con RMSLE). Los análisis publicados se centran en construir
atributos nuevos, como la distancia entre origen y destino o la hora y el día de
la semana, y en detectar datos anómalos~\cite{eda_blog}. Otros combinan la
generación automática de atributos con modelos de \textit{boosting}~\cite{alteryx_taxi}.
Algo parecido, aunque con trayectorias GPS de taxis de Oporto, se planteó en el
reto ECML/PKDD 2015~\cite{lam2015}. Nuestro enfoque es distinto: en lugar de
predecir la duración exacta, solo queremos saber si el viaje será corto, medio o
largo.

\subsection{Descripción del conjunto de datos}
El conjunto de datos original cuenta con 1\,458\,644 registros. Tras una inspección inicial, no se detectan valores nulos explícitos en ninguna de las columnas. Sin embargo, se observan anomalías e inconsistencias que requieren tratamiento:
\begin{itemize}
    \item \textbf{Atributos temporales:} Abarcan desde el 1 de enero hasta el 30 de junio de 2016. La variable \texttt{dropoff\_datetime} representa el final del viaje. Dado que en un entorno de producción real esta información no está disponible al iniciar la carrera, debe eliminarse para evitar fuga de información (\textit{data leakage}).
    \item \textbf{Duración de trayectos (\texttt{trip\_duration}):} Variable numérica continua con una media de 959 segundos. Presenta valores atípicos extremos, con viajes de 1 segundo (ruido o cancelaciones) y máximos de varios días (outliers).
    \item \textbf{Coordenadas geográficas:} Existen registros con longitudes o latitudes que caen fuera de los límites lógicos de Nueva York e incluso del océano.
    \item \textbf{Pasajeros (\texttt{passenger\_count}):} Se registran viajes con 0 pasajeros, lo cual representa un registro no válido en el contexto de un servicio de taxi.
\end{itemize}

\subsection{Examinar correlaciones en los datos}
El análisis exploratorio previo revela ciertas dependencias clave que guiarán la selección de atributos:
\begin{enumerate}
    \item \textbf{Distancia y Duración:} Existe una fuerte dependencia entre la distancia del trayecto y la clase de duración. Calcular la distancia geográfica real (por ejemplo, mediante la fórmula de Haversine) aportará alto valor predictivo.
    \item \textbf{Variables temporales:} Las horas punta (17:00 a 19:00 horas) aumentan la probabilidad de pertenecer a las clases \textit{medium} o \textit{long} incluso para distancias cortas debido a la congestión urbana. Asimismo, el día de la semana influye en los patrones de tráfico.
\end{enumerate}

\section{Preparación de Datos}
\subsection{Creación de conjuntos de entrenamiento y prueba}
Para asegurar una evaluación rigurosa y evitar el sesgo, se ha realizado una partición aleatoria del 80\,\% para entrenamiento y 20\,\% para prueba (\texttt{seed = 42L}). Esta división se ejecuta inmediatamente tras descartar posibles filas nulas completas y asignar la clase objetivo, pero \textbf{antes} de aplicar filtros profundos o transformaciones.

\subsection{Limpieza de datos}
Se ha aplicado un filtrado para manejar el ruido, las inconsistencias y los outliers detectados en la exploración:
\begin{itemize}
    \item \textbf{Inconsistencias:} Se descartan viajes con 0 pasajeros (\texttt{passenger\_count > 0}).
    \item \textbf{Outliers espaciales:} Se acotan las coordenadas a un recuadro válido para el área de Nueva York ($\text{Lat} \in [40.4, 40.9]$ y $\text{Lon} \in [-74.3, -73.7]$).
    \item \textbf{Outliers temporales:} Se limitan los viajes a duraciones lógicas, acotando entre 10 segundos y 24 horas (86\,400 segundos).
\end{itemize}

Aplicando estos mismos filtros al conjunto de prueba, obtenemos las métricas de seguimiento requeridas: partiendo de 291\,657 instancias iniciales en el test, se han eliminado 744 registros inconsistentes y atípicos. Por tanto, $NTAIE + NTOE = 744$.
\[
\text{Tasa de No Clasificados} = \frac{744}{291\,657} \approx 0,00255 \text{ (0,25\,\%)}
\]
Esta tasa del 0,25\,\% es completamente marginal, cumpliendo sobradamente el criterio exigido de ser al menos un orden de magnitud inferior a la tasa de error esperada del modelo.

\subsection{Selección de atributos}
Se descartan los atributos \texttt{id} (sin valor predictivo) y \texttt{dropoff\_datetime} (evita \textit{data leakage}). Para capturar las correlaciones analizadas, se generan nuevos atributos sintéticos a partir de los datos originales:
\begin{itemize}
    \item \textbf{Espaciales:} \texttt{distance\_km} (distancia en línea recta calculada mediante la fórmula de Haversine empleando las coordenadas de origen y destino).
    \item \textbf{Temporales:} \texttt{pickup\_hour}, \texttt{pickup\_dayofweek} y \texttt{pickup\_month} extraídos de la fecha de recogida.
\end{itemize}

\subsection{Transformación de datos}
Se emplean las abstracciones de \textit{Transformers} y \textit{Estimators} proporcionadas por Spark ML para adaptar los datos al formato numérico-vectorial requerido:
\begin{enumerate}
    \item \texttt{StringIndexer}: Utilizado para convertir la variable objetivo nominal en un índice numérico (\texttt{label}) y para codificar la variable booleana \texttt{store\_and\_fwd\_flag}.
    \item \texttt{VectorAssembler}: Empleado para fusionar todas las características predictoras discretas y continuas en un único vector denso denominado \texttt{features}.
\end{enumerate}
En cumplimiento estricto del marco metodológico, los \textit{Estimators} (\texttt{StringIndexer}) se han ajustado (\texttt{fit}) \textbf{exclusivamente} sobre el conjunto de entrenamiento, aplicando posteriormente sus reglas de transformación (\texttt{transform}) de manera idéntica tanto al subconjunto de entrenamiento como al de prueba.

\section{Conclusiones}
En esta primera etapa se ha completado satisfactoriamente el ciclo de preparación de datos en Spark. Tras una exploración detallada, se ha llevado a cabo una limpieza rigurosa manejando valores inconsistentes, outliers espaciales y valores temporales extremos. La correcta separación en entrenamiento y prueba, unida a la ingeniería de características (cálculo de distancias y desglose temporal) y la vectorización final de atributos, deja los conjuntos listos y estructurados en formato \texttt{features}/\texttt{labels} para alimentar de forma segura los algoritmos de clasificación en la siguiente fase del proyecto, manteniendo una pérdida de información en prueba (tasa de no clasificados) del 0,25\,\%.

\section*{Dedicación del Grupo}
A continuación se detalla la contribución equitativa de los integrantes del equipo:

\begin{table}[H]
\centering
\small
\begin{tabular}{lcp{7.5cm}}
\toprule
\textbf{Estudiante} & \textbf{Horas / \%} & \textbf{Tareas Principales Realizadas} \\
\midrule
Gonzalo Agúndez Sáez & 18 h (33,3\,\%) & Desarrollo de scripts de preprocesado en Spark (Scala) y cálculo de métricas de limpieza ($NTAIE$/$NTOE$). \\
\addlinespace
Paula Negro Asegurado & 18 h (33,3\,\%) & Ingeniería de atributos sintéticos, maquetación integral del documento LaTeX y resumen ejecutivo. \\
\addlinespace
Diego Sanz Santos & 18 h (33,3\,\%) & Exploración inicial de datos en entorno Python, análisis de distribuciones y extracción de correlaciones. \\
\bottomrule
\end{tabular}
\caption{Desglose de dedicación y carga de trabajo por integrante.}
\end{table}


\clearpage
\bibliographystyle{unsrt}
\bibliography{referencias}

\clearpage
\appendix
\section{Apéndices}
% Ideal para métricas de limpieza como NTAIE y NTOE
\end{document}